% ---- score2 standard preamble (from: python3 score2_tools.py preamble) ----
\documentclass[10pt]{extarticle}
\usepackage[a4paper,margin=1.5cm]{geometry}
\usepackage[T1]{fontenc}
\usepackage{amsmath,amssymb}
\usepackage{graphicx}
\usepackage[export]{adjustbox}
\usepackage{multicol}
\usepackage{enumitem}
\usepackage{booktabs,array,tabularx}
\setlength{\columnsep}{0.7cm}
\setlength{\parindent}{0pt}
\setlength{\parskip}{3pt}
\pagestyle{empty}
\makeatletter
% figure registry (figdefs.tex is generated by the extract step)
\newcommand{\qfigdef}[4]{\@namedef{qf@file@#1}{#2}\@namedef{qf@w@#1}{#3}\@namedef{qf@wide@#1}{#4}}
\newcommand{\qf@check}[1]{\@ifundefined{qf@file@#1}{\PackageError{score2}{Unknown figure id #1}{}}{}}
\newcommand{\qf@img}[4][]{\includegraphics[width=#2\dimexpr\@nameuse{qf@w@#4}pt\relax,max width=#3,max height=0.45\textheight,#1]{\@nameuse{qf@file@#4}}}
% Full-width questions: finalize wraps every question that contains a wide figure in \qfullbegin...\qfullend,
% so the WHOLE question (stem, figures, options) is set across the page in reading order, then the two columns
% resume. The question's head - its number, stem and first figure (\qfullsplit marks the end of the head) -
% is kept on one page, so a number is never left at a page foot with its figure overleaf; the rest may break
% across pages normally, so no big gaps are left. Without a \qfullsplit, the whole question is kept together
% if it is under 45% of a page. Never write these yourself.
\newif\ifqf@full
\newif\ifqf@boxed
\newsavebox{\qf@box}
\newcommand{\qf@open}{\global\qf@fulltrue\par\global\qf@boxedtrue
  \setbox\qf@box\vbox\bgroup\hsize\textwidth\linewidth\textwidth}
\newcommand{\qf@emit}[1]{\par\egroup\global\qf@boxedfalse
  \ifdim\dimexpr\ht\qf@box+\dp\qf@box\relax<#1\textheight\noindent\usebox{\qf@box}\par\else\unvbox\qf@box\fi}
\newcommand{\qfullsplit}{\ifqf@boxed\qf@emit{0.85}\fi}
\newcommand{\qf@close}{\ifqf@boxed\qf@emit{0.45}\fi\global\qf@fullfalse}
\newcommand{\qfullbegin}{\end{multicols}\qf@open}
\newcommand{\qfullend}{\qf@close\begin{multicols}{2}}
\newcommand{\qfullstart}{\qf@open}
\newcommand{\qfulldone}{\qf@close}
% \qfig[scale]{id}: block figure, centred, directly where it is written (size set by the tool).
% Fallback if a wide figure is not inside a full-width question: break out of the columns in place - no floats.
\newcommand{\qfig}[2][1]{\qf@check{#2}%
  \ifqf@full\qf@blk{#1}{#2}\else
  \ifnum\@nameuse{qf@wide@#2}=1 \expandafter\@firstoftwo\else\expandafter\@secondoftwo\fi
  {\end{multicols}\begin{center}\qf@img{#1}{\textwidth}{#2}\end{center}\begin{multicols}{2}}%
  {\qf@blk{#1}{#2}}\fi}
\newcommand{\qf@blk}[2]{\par\smallskip{\centering\qf@img{#1}{\linewidth}{#2}\par}\smallskip}
% \qfigi[scale]{id}: inline figure (inside option lists / tables / minipages); never breaks columns.
\newcommand{\qfigi}[2][1]{\qf@check{#2}\qf@img[valign=c]{#1}{\linewidth}{#2}}
\makeatother
% structure macros - lint reads these to check completeness and figure placement
\newcommand{\papersec}[1]{\par\bigskip{\centering\Large\bfseries PAPER-#1\par}\medskip}
\newcommand{\partsec}[2]{\par\medskip\noindent{\large\bfseries PART-#1: #2}\par\smallskip}
\newcommand{\sect}[1]{\par\smallskip\noindent\textbf{#1}\par}
\newcommand{\qn}[1]{\par\smallskip\noindent\textbf{#1.}~}
\newcommand{\paraq}[2]{\par\smallskip\noindent\textbf{Paragraph for Questions #1 and #2}\par}
\newenvironment{qoptions}{\begin{itemize}[leftmargin=1.8em,itemsep=1pt,topsep=2pt,label={}]}{\end{itemize}}
\newcommand{\opt}[2]{\item[(#1)] #2}
% figure registry (generated by score2_tools.py finalize)
\qfigdef{p1_phy_s1_q01_1}{Pre-Score (2023-24) - Test 02 (15 Oct 2023)-img/p1_phy_s1_q01_1.png}{129.4}{0}
\qfigdef{p1_phy_s1_q02_1}{Pre-Score (2023-24) - Test 02 (15 Oct 2023)-img/p1_phy_s1_q02_1.png}{113.6}{0}
\qfigdef{p1_phy_s1_q04_1}{Pre-Score (2023-24) - Test 02 (15 Oct 2023)-img/p1_phy_s1_q04_1.png}{55.6}{0}
\qfigdef{p1_phy_s1_q06_1}{Pre-Score (2023-24) - Test 02 (15 Oct 2023)-img/p1_phy_s1_q06_1.png}{182.2}{0}
\qfigdef{p1_phy_s2_q07_1}{Pre-Score (2023-24) - Test 02 (15 Oct 2023)-img/p1_phy_s2_q07_1.png}{119.0}{0}
\qfigdef{p1_phy_s2_q08_1}{Pre-Score (2023-24) - Test 02 (15 Oct 2023)-img/p1_phy_s2_q08_1.png}{116.3}{0}
\qfigdef{p1_phy_s3_q05_1}{Pre-Score (2023-24) - Test 02 (15 Oct 2023)-img/p1_phy_s3_q05_1.png}{200.8}{0}
\qfigdef{p1_chem_s2_q07_1}{Pre-Score (2023-24) - Test 02 (15 Oct 2023)-img/p1_chem_s2_q07_1.png}{306.4}{1}
\qfigdef{p2_phy_s1_q04_1}{Pre-Score (2023-24) - Test 02 (15 Oct 2023)-img/p2_phy_s1_q04_1.png}{171.7}{0}
\qfigdef{p2_phy_s2_q09_1}{Pre-Score (2023-24) - Test 02 (15 Oct 2023)-img/p2_phy_s2_q09_1.png}{107.3}{0}
\qfigdef{p2_chem_s3_q02_1}{Pre-Score (2023-24) - Test 02 (15 Oct 2023)-img/p2_chem_s3_q02_1.png}{171.4}{0}
% ---- body skeleton ----
% \begin{document}
% \begin{center}{\Large\bfseries JEE (Main+Advanced) Score-II --- Test NN (DD Mon YYYY)}\end{center}
% \begin{multicols}{2}
%   \papersec{1}   (only for two-paper tests)
%   \partsec{1}{PHYSICS} \sect{Section-I (i): Single Correct} \qn{1} ... \qfig{p1_phy_q01_1} ...
% \end{multicols}
% \end{document}

\begin{document}
\begin{center}{\Large\bfseries Pre-Score (2023-24) --- Test 02 (15 Oct 2023)}\end{center}
\begin{multicols}{2}
\papersec{1}
\partsec{1}{PHYSICS}
\sect{SECTION-I(i)}
\qn{1} Consider the situation shown in figure source of sound `S' and listener `O' both are moving towards a wall with velocity $v_S$ and $v_O$ respectively. The wall is moving away from them with velocity $u$. The wind is blowing with a velocity $v_W$ towards the wall. The frequency of sound emitted by source when it is stationary is $f_O$. If the velocity of sound relative to air is $v$, then
\qfig{p1_phy_s1_q01_1}
\begin{qoptions}
\opt{A}{frequency of sound received by listener after reflection from wall is $\Bigl(\dfrac{v-v_W+v_O}{v+v_W-v_S}\Bigr)\Bigl(\dfrac{v+v_W-u}{v-v_W+u}\Bigr)f_O$}
\opt{B}{Wavelength of sound received by wall is $\dfrac{(v-v_S)}{f_O}$}
\opt{C}{Wavelength of sound received by listener after reflection from wall is $\dfrac{(v-v_W+u)(v+v_W-v_S)}{(v+v_W-u)f_O}$.}
\opt{D}{Frequency of sound received by listener after reflection from wall $\dfrac{(v+v_W+u)(v-v_W+v_S)}{(v+v_W-u)}f_O$}
\end{qoptions}
\qn{2} A metal rod of length 135 cm is clamped at two points as shown in figure. One clamp is at an end and other is at 45 cm from other end. The Young's modulus of elasticity and density of the rod are $Y=1.6\times10^{11}$ N/m$^2$ and $\rho=2.5$ gm/cc respectively.
\qfig{p1_phy_s1_q02_1}
\begin{qoptions}
\opt{A}{the rod will resonate if it is given longitudinal vibrations of frequency 48.84 kHz.}
\opt{B}{the rod will resonate if it is given longitudinal vibrations of frequency 26.64 kHz.}
\opt{C}{during resonance the number of complete loops between clamps may be 6.}
\opt{D}{the maximum wavelength corresponding to vibrations of the rod is 90 cm.}
\end{qoptions}
\qn{3} Twins Anna and Tom measure each other's heights each birthday. When they were 12 Anna's height was measured as 151.5 cm and Tom's was measured as 151 cm. On their 16th birthday Anna's height is measured as 167.5 cm and Tom's is measured as 168.5 cm. The uncertainty in the height measurements is 5 mm. Anna and Tom have always been very competitive about their heights, so Tom crows with glee after being measured, saying ``Finally, I'm taller than you! We were always about the same height, but now I'm taller than you!'' Anna responds, ``You're crazy! I was clearly taller than you when we were 12.'' Which of the following statements is most correct?
\begin{qoptions}
\opt{A}{Tom is probably right and Anna is wrong, they were about the same height at 12, with a difference of 0.5 cm $\pm$ 1.0 cm. Now the measurements show the difference to be 1.0 cm $\pm$ 1.0 cm, more accurate measurements would resolve the dispute over who is taller now.}
\opt{B}{Tom is probably right and Anna is wrong, they were about the same height at 12, with a difference of 0.5 cm $\pm$ 1.0 cm. Now the measurements show the difference to be 1.0 cm $\pm$ 1.0 cm, however, taking more measurements, regardless of how accurate, couldn't help resolve the dispute over who is taller now.}
\opt{C}{Anna is right and Tom is wrong, Anna was taller at 12, with a difference of 0.5 cm $\pm$ 0.5 cm. Now Tom is taller by 1.0 cm $\pm$ 0.5 cm.}
\opt{D}{Anna is right and Tom is wrong, she was 0.5 cm taller than Tom when they were 12. Tom is now 1.0 cm taller than Anna.}
\end{qoptions}
\qn{4} A soap film, with constant refractive index $\mu=1.5$, is contained within a vertical rectangular wire frame and has a wedge shaped profile, small wedge angle $\phi$, figure. The film drains under gravity, being thicker at the bottom than at the top, where the thickness can be considered to be zero. The film is viewed, by reflection, with red and violet light at near normal incidence; the angle of incidence being exceedingly small. Interference results from light reflected from the outer surface and that emerging after reflection at the inner surface of the film. The first violet, wavelength 420 nm, constructive interference band is observed at a distance $x_v=3.0$ cm from the top edge of the film. The wavelength of the red light is 700 nm.
\qfig{p1_phy_s1_q04_1}
\begin{qoptions}
\opt{A}{The thickness of the film at the lowest order constructive violet interference is 70 nm.}
\opt{B}{The film wedge angle $\phi$ is $2.33\times10^{-6}$ radian.}
\opt{C}{For the red light, the distance of lowest order constructive interference from the top of the film is 2.5 cm.}
\opt{D}{The soap film drains resulting in a decrease in wedge angle at a rate of $8\times10^{-5}$ rad/min. The initial speed of observed red interference fringes is 5.7 cm/s.}
\end{qoptions}
\qn{5} For the data given in the table below.
\begin{center}\small
\begin{tabular}{|l|c|c|}
\hline medium & \begin{tabular}[c]{@{}c@{}}density\\ kg m$^{-3}$\end{tabular} & \begin{tabular}[c]{@{}c@{}}sound speed\\ ms$^{-1}$\end{tabular} \\\hline
air & 1.2 & 300 \\\hline
concrete & 2400 & 3600 \\\hline
water & 1000 & 1500 \\\hline
\end{tabular}
\end{center}
\begin{qoptions}
\opt{A}{If sound is incident from concrete to air, then the percentage of energy reflected at the boundary is approx. 80\%.}
\opt{B}{If sound is incident from air to concrete, then the percentage of energy reflected at the boundary is approx. 70\%.}
\opt{C}{If sound is incident from concrete to water, then the percentage of energy reflected at the boundary is approx. 50\%.}
\opt{D}{If sound is incident from concrete to water, then the percentage of energy reflected at the boundary is approx. 40\%.}
\end{qoptions}
\qn{6} A vibrating string 50.0 cm long is under a tension of 1.00 N. The results from five successive stroboscopic pictures are shown in figure. The strobe rate is set at 1000 flashes per minute, and observations reveal that the maximum displacement occurred at flashes 1 and 5 with no other maxima in between.
\qfig{p1_phy_s1_q06_1}
\begin{qoptions}
\opt{A}{The time period of oscillation of string is 0.24 sec.}
\opt{B}{The speed of travelling waves on the string is $\dfrac{25}{24}$ m/s.}
\opt{C}{The mass of the string is 0.92 kg.}
\opt{D}{The point $P$ at anti-node moves with a speed of $6.25\pi$ cm/sec when the string is at position-3.}
\end{qoptions}
\sect{SECTION-I(ii)}
\paraq{7}{8}
The spherometer is used to measure thickness of thin plates and radius of curvature of spherical surfaces. The principle of measurement in the spherometer is same as that of the screw gauge. The circular scale of the spherometer is rotated by a right-handed screw. When the circular scale rotates in the counter-clockwise direction it moves up. Before making any measurement, the three legs of the circular scale are placed on a plane level surface and the screw is rotated so that its tip touches the surface. In this position, if the zero mark of the main scale does not coincide with the zero mark of the circular scale, the instrument is said to possess a zero error. If the zero mark of the circular scale lies a little ahead (in the direction of rotation) of that of the main scale, the zero error is taken to be positive. If the zero mark of the circular scale lies a little behind that of the main scale, the zero error is negative. In the spherometer shown in fig(a), the total number of divisions on the circular scale of the spherometer is 200 and 5 divisions of the main scale is equal to 10 mm. When the spherometer is given one rotation it travels a distance of one division on the main scale. This is called the pitch of the spherometer. In the level position of the legs and the screw, the zero mark of the circular scale lies 50 divisions (of the circular scale) behind the zero mark of the main scale. When thickness of the plate is measured, 45$^{\mathrm{th}}$ division of the circular scale lies in between 2$^{\mathrm{nd}}$ and 3$^{\mathrm{rd}}$ division of the main scale.
\qn{7} The thickness of the plate in figure (a) is
\qfig{p1_phy_s2_q07_1}
\begin{qoptions}
\opt{A}{4.45 mm}
\opt{B}{4.95 mm}
\opt{C}{3.95 mm}
\opt{D}{none of these}
\end{qoptions}
\qn{8} Another spherometer shows the scale, the screw and all the three legs touch the plane surface as shown in figure (b). The least count of the spherometer is given to be 0.01 mm. Then zero error in the instrument is
\qfig{p1_phy_s2_q08_1}
\begin{qoptions}
\opt{A}{0.98 mm}
\opt{B}{$-0.98$ mm}
\opt{C}{$-0.02$ mm}
\opt{D}{0.02 mm}
\end{qoptions}
\paraq{9}{10}
A colloid consists of particles of one type of substance dispersed in another substance. Suspensions of electrically charged microspheres (microscopic spheres, such as polystyrene) in a liquid such as water can form a colloidal crystal when the microspheres arrange themselves in a regular repeating pattern under the influence of the electrostatic force. Colloidal crystals can selectively manipulate different wavelengths of visible light. Just as we can study crystal-line solids by using Bragg reflection of x rays, we can study colloidal crystals through Bragg scattering of visible light from the regular arrangement of charged microspheres. Because the light is traveling through a liquid when it experiences the path differences that lead to constructive interference, it is the wavelength in the liquid that determines the angles at which Bragg reflections are seen. In one experiment, laser light with a wavelength in vacuum of 650 nm is passed through a sample of charged polystyrene spheres in water. A strong interference maximum is then observed when the incident and reflected beams make an angle of $39^\circ$ with the colloidal crystal planes.
\qn{9} Why is visible light, which has much longer wavelengths than x rays do, used for Bragg reflection experiments on colloidal crystals?
\begin{qoptions}
\opt{A}{The microspheres are suspended in a liquid, and it is more difficult for x rays to penetrate liquid than it is for visible light.}
\opt{B}{The irregular spacing of the microspheres allows the longer wavelength visible light to produce more destructive interference than can x rays.}
\opt{C}{The microspheres are much larger than atoms in a crystalline solid, and in order to get interference maxima at reasonably large angles, the wavelength must be much longer than the size of the individual scatterers.}
\opt{D}{The microspheres are spaced more widely than atoms in a crystalline solid, and in order to get interference maxima at reasonably large angles, the wavelength must be comparable to the spacing between scattering planes.}
\end{qoptions}
\qn{10} When the light is passed through the bottom of the sample container, the interference maximum is observed to be at $41^\circ$; when it is passed through the top, the corresponding maximum is at $37^\circ$. What is the best explanation for this observation?
\begin{qoptions}
\opt{A}{The microspheres are more tightly packed at the bottom, because they tend to settle in the suspension.}
\opt{B}{The microspheres are more tightly packed at the top, because they tend to float to the top of the suspension.}
\opt{C}{The increased pressure at the bottom makes the microspheres smaller there.}
\opt{D}{The maximum at the bottom corresponds to $m=2$, whereas the maximum at the top corresponds to $m=1$.}
\end{qoptions}
\sect{SECTION-II}
\qn{1} Two tuning forks of frequency 250 Hz and 256 Hz are sounded together. What shall be the time interval (in sec.) between maximum and minimum sound at a point?
\qn{2} A bicycle pump of cross-sectional area 4.0 cm$^2$ has one end sawn off and a cork is fitted into the end. The piston is pushed slowly inwards and the cork is fired out with a popping sound which has a fundamental frequency of 512 Hz. The initial distance between the cork and the piston is 25 cm, with atmospheric pressure equal to $1.0\times10^5$ Pa and the speed of sound in air being 330 ms$^{-1}$. The force required to eject the cork (in N) is
\qn{3} Monochromatic light of wavelength 600 nm is incident on two vertical slits hence producing two coherent sources. Before the light leaving these slits overlaps and interferes, each beam passes through a tube 5.0 cm long. One of the tubes is now gradually evacuated and it is noted that the fringe pattern shifts 25 fringes. The refractive index of air is more than vaccum by $\alpha\times10^{-3}$. The value of $\alpha$ is
\qn{4} Two whistles, both of frequency $f=2.00$ kHz, are situated 3.00 m apart and blown simultaneously. An observer travelling along a line parallel to the line joining the whistles, and at a distance approximately 20.0 m opposite the whistles, detects minima in sound intensity at a series of points spaced 1.14 m apart. Calculate the speed (in m/s) of sound in air.
\qn{5} Determining the wind speed in a tornado using conventional meteorological instruments is difficult (since the tornado is small in size and moves) and unsafe. It is proposed to measure it from afar using a portable radar, since there is a lot of dust and small objects inside the tornado that reflects ultrasonic waves. Radar emits ultrasonic waves at a frequency $f_o=30$ kHz. The spectrum of the signal reflected from the tornado is shown in the figure. Find the maximum wind speed (in m/s) in the tornado. (Take : speed of sound $=320$ m/s)
\qfig{p1_phy_s3_q05_1}
\qn{6} The vernier of a vernier scale is divided into 10 divisions which coincide with 9 division of the main scale, each main scale division being 0.5 mm. When the two jaws of the instrument are in contact with each other, the 4th division of the vernier scale coincides with a main scale division and the zero of the vernier lies to the right of the zero of the main scale. When a sphere is inserted between the jaws, the zero of vernier scale lies slightly to the right of the 1.8 cm and the sixth vernier division coincides with a main division. The diameter (in cm) of sphere will be :-
\qn{7} The objective of telescope has diameter 12.2 cm. The distance (in km) at which two small green object placed 30cm apart can be barely resolved by telescope, assuming the resolution to be limited by diffraction by objective only ($\lambda=5.4\times10^{-5}$ cm)
\qn{8} An electromagnetic standing wave in a certain material has frequency $1.20\times10^{10}$ Hz and speed of propagation $2.10\times10^{8}$ m/s. The distance between a nodal plane of $\vec E$ and the closest nodal plane of $\vec B$ (in cm) is :
\partsec{2}{CHEMISTRY}
\sect{SECTION-I(i)}
\qn{1} When Pb$_{10}$O$_3$(CO$_3$)$_6$(OH)$_2$ react with dil. HNO$_3$ at room temperature it produces :-
\begin{qoptions}
\opt{A}{NO$_2$}
\opt{B}{PbO$_2$}
\opt{C}{Pb(NO$_3$)$_2$}
\opt{D}{CO$_2$}
\end{qoptions}
\qn{2} A solid mixture of silver nitrate and an alkali metal nitrate was dissolve in a small volume of water when an excess of an alkali metal salt (Q) was added to the solution, a silver salt precipitated from solution. After the precipitate was collected, dried and weighted, its mass was equal to the mass of the initial solid mixture.\\
Which of the following substance could be the alkali metal salt (Q) added?
\begin{qoptions}
\opt{A}{Potassium fluoride}
\opt{B}{Sodium cynide}
\opt{C}{Potassium bromide}
\opt{D}{Rubidium iodide}
\end{qoptions}
\qn{3} The compound(s) having [Xe]4f$^1$ configuration on metal ion is/are\\
(Atomic number of Th = 90, Ce = 58, Lu = 71, Pr = 59) :-
\begin{qoptions}
\opt{A}{Na$_3$[Ce(NO$_3$)$_6$]}
\opt{B}{Na$_3$[LuCl$_6$]}
\opt{C}{PrO$_2$}
\opt{D}{(NH$_4$)$_4$[ThF$_8$]}
\end{qoptions}
\qn{4} Metal Salt$_{\mathrm{(aq.)}}$ + H$_2$S/acidic solution $\longrightarrow$ [Black ppt.] $\xrightarrow[\mathrm{dil.\ HNO_3}]{\Delta\ \mathrm{with}}$ Black ppt disappear\\
Metal Salt$_{\mathrm{(aq.)}}$ can be :-
\begin{qoptions}
\opt{A}{Pb(CH$_3$COO)$_2$}
\opt{B}{NiCl$_2$}
\opt{C}{CuSO$_4$}
\opt{D}{Hg$_2$(NO$_3$)$_2$}
\end{qoptions}
\qn{5} Diatomic paramagnetic gas is produce by :-
\begin{qoptions}
\opt{A}{Reaction of NH$_3$ \& oxygen using Pt/Rh catalyst}
\opt{B}{Thermal decomposition of KClO$_3$}
\opt{C}{Heating ammonium nitrite}
\opt{D}{Dissolving gold in aqua regia}
\end{qoptions}
\qn{6} The electrons in \underline{\hspace{1.5cm}} are too strongly bound to get excited by flame.\\
Select correct to fill the blank space.
\begin{qoptions}
\opt{A}{Lithium}
\opt{B}{Beryllium}
\opt{C}{Barium}
\opt{D}{Magnesium}
\end{qoptions}
\sect{SECTION-I(ii)}
\paraq{7}{8}
The substituted Halosilanes undergoes hydrolysis to corresponding substituted Silanols, which on condensation produce silicones. The properties of these silicones depends upon the degree of polymerization and stereo chemistry more --OH groups in the Silanols naturally lead to condensation in one or more direction.
\qfullbegin
\qn{7} Which polymeric substance not required Mono Alkyl/Aryl Silanols during its formation :-\\
(P) Silicon Resins \hfill (Q) Cross link silicone
\qfig{p1_chem_s2_q07_1}
\qfullsplit
Select most appropriate option :-
\begin{qoptions}
\opt{A}{(P) \& (Q)}
\opt{B}{(R) \& (S)}
\opt{C}{(Q), (R) \& (S)}
\opt{D}{(P) \& (R)}
\end{qoptions}
\qfullend
\qn{8} Which substance has structural resemblance with ``Pyrosilicate'' ion :-
\begin{qoptions}
\opt{A}{Polydimethylsiloxane silicone polymer}
\opt{B}{Condensation product of trimethylsilanols}
\opt{C}{Condensation product of monoalkyl silanols}
\opt{D}{(Me$_2$SiO)$_n$ (where $n=3$)}
\end{qoptions}
\paraq{9}{10}
Coloured solution (S$_1$) contain two unknown metal ions treated with dil. HCl produces no ppt., Coloured solution (S$_1$) react with H$_2$S in acidic medium both ions are precipitated, precipitate of both ions are soluble in dil. HNO$_3$ to gives coloured solution (S$_2$) and turbidity of sulphur.\\
Solution (S$_2$) produces deep blue solution (S$_3$) \& white ppt. (P$_1$) with excess of NH$_4$OH. After filtration ppt. (P$_1$) dissolve in dil. H$_2$SO$_4$ \& produce black ppt. (P$_2$) with sodium tetrahydroxidostanate(II).\\
Solution (S$_3$) first treated with excess dil. H$_2$SO$_4$ and then with potassium hexacynidoferrate(II) produce brown ppt. (P$_3$).
\qn{9} Dark Blue solution (S$_3$) is :-
\begin{qoptions}
\opt{A}{[Ni(NH$_3$)$_6$](NO$_3$)$_2$}
\opt{B}{Hg[Co(SCN)$_4$]}
\opt{C}{CrO$_5$}
\opt{D}{[Cu(NH$_3$)$_4$](NO$_3$)$_2$}
\end{qoptions}
\qn{10} Oxidation number of unknown metal ions in precipitate (P$_2$) and (P$_3$) respectively :-
\begin{qoptions}
\opt{A}{0, $+2$}
\opt{B}{$+3$, 0}
\opt{C}{$+4$, $+2$}
\opt{D}{$+5$, $+3$}
\end{qoptions}
\sect{SECTION-II}
\qn{1} The value of ``$n$'' in P$_{(n)}$O$_{18}^{-6}$ is \underline{\hspace{1cm}}.\\
The value of ``$m$'' in S$_{(m)}$O$_9$ is \underline{\hspace{1cm}}.\\
Sum of $n$ \& $m$ is your answer.
\qn{2} Consider given reaction from metal M\\
$\underset{\text{(A 3d series metal)}}{\mathrm{M}}+2\mathrm{HCl_{(aq.)}}\ \text{(dil.)}\xrightarrow{\text{Air excluded}}\underset{\text{(Blue)}}{\mathrm{M^{+2}_{(aq.)}}}+\mathrm{H_2}\uparrow$\\
$\mathrm{M^{+2}_{(aq.)}}+2\mathrm{H^+_{(aq.)}}\xrightarrow{\text{slowly}}\mathrm{M^{+3}_{(aq.)}}+\mathrm{H_2}\uparrow$\\
$\mathrm{M^{+3}_{(aq.)}}\xrightarrow{\mathrm{NaBO_3+NaOH}}$ Yellow solution (P)\\
Find sum of the number of d-electrons in $\mathrm{M^{+3}_{(aq.)}}$ \& oxidation number of metal in solution (P).
\qn{3} 100 gram sample of table salt sample which contain 15\% of potassium iodate is dissolve in 1 litre of dil. HCl and treated with excess KI.\\
Calculate amount of H$_2$O produce in above reaction of iodine containing salts.\\
(given that atomic mass of K = 39, I = 127, O = 16, Cl = 35.5, H = 1)
\qn{4} (E)$_{(s)}$ (Non metallic element which has octa -- atomic molecules with puckered ring) $\xrightarrow{\mathrm{KOH_{(aq.)}}}$ (P)$_{(aq.)}$ salt + (Q)$_{(aq.)}$ salt\\
(P) + dil. HCl $\to$ Gas--(R) + (E) $\downarrow$\\
(Q) + dil. HCl $\to$ Gas--(S)\\
(R) + MnO$_4^-$/H$^+$ $\to$ (T) + Mn contain product\\
(P) + NaI$_{3\mathrm{(aq.)}}$ $\xrightarrow{\mathrm{H^+}}$ (U) + Iodine contain product\\
Find sum of non metal atom (E) -- Oxygen linkages (i.e. E--O linkages) in (P), (Q), (R), (S), (T) and (U)
\qn{5} $\bullet$ Metal (M$_1$) of 3d series form unstable MO$_4^{2-}$ in alkaline medium but readily decompose to M$_2$O$_3$ \& O$_2$\\
$\bullet$ Metal (M$_2$) of 3d series has half filled t$_{2g}$ level in $[\mathrm{M(H_2O)_6}]^{+2}$ ion\\
$\bullet$ Metal (M$_3$) of 3d series produce stable $\mathrm{M^+_{(aq.)}}$ ion rather than $\mathrm{M^+_{(aq.)}}$ and $\mathrm{M^+_{(aq.)}}$ is disproportionate to M and $\mathrm{M^{+2}_{(aq.)}}$.\\
Identify metal M$_1$, M$_2$ and M$_3$ and write the sum of atomic number of M$_1$, M$_2$ and M$_3$.
\qn{6} HNO$_3$ $\xrightarrow{\mathrm{P_4O_{10}}}$ Ionic solid (X) + other product - - -\\
Ionic solid (X) $\xrightarrow{\text{Na metal}}$ Ionic solid(Y) + other product - - -\\
Write the sum of the formula weight of ionic solid (X) \& ionic solid (Y).\\
(Given atomic mass of Na = 23, H = 1, N = 14, O = 16, P = 31)
\qn{7} Halogen ``X$_2$'' + Halogen ``Y$_2$'' (Excess) $\longrightarrow$ Interhalogen compound (Z)\\
Z is a orange solid which is dimerises as ``Y''-bridged dimer.\\
Find total number of lone pairs in dimer of ``Z''.
\qn{8} 7.45 gm of KCl (crystals) are heated with excess MnO(OH)$_2$ and H$_2$SO$_{4\mathrm{(aq.)}}$. The gas so produced was then passed through a solution of KI$_{\mathrm{(aq.)}}$. Calculate weight of non polar diatomic solid substance can be produce in final step.\\
(Given that atomic mass of K = 39, Cl = 35.5, H = 1, O = 16, S = 32, I = 127, Mn = 55)
\partsec{3}{MATHEMATICS}
\sect{SECTION-I(i)}
\qn{1} If roots of the equation $2x^3+4x^2-5x+8=0$ ae $\alpha,\beta,\gamma$ then
\begin{qoptions}
\opt{A}{$\sqrt{(2-\alpha)(2-\beta)(\gamma-2)(1-\alpha^2)(1-\beta^2)(1-\gamma^2)}=22.5$}
\opt{B}{$1+\alpha,1+\beta,1+\gamma$ are roots of $2x^3-2x^2-7x+15=0$}
\opt{C}{$\sqrt{(2-\alpha)(2-\beta)(2-\gamma)(1-\alpha^2)(1-\beta^2)(1-\gamma^2)}=22.5$}
\opt{D}{$1+\alpha,1+\beta,1+\gamma$ are roots of $2x^3-6x^2-7x-15=0$}
\end{qoptions}
\qn{2} If the value of $\Bigl(1+\dfrac23+\dfrac6{3^2}+\dfrac{10}{3^3}+\cdots\text{ upto }\infty\Bigr)^{\log_{(0.25)}\left(\frac13+\frac1{3^2}+\frac1{3^3}+\cdots\text{ upto }\infty\right)}$ is $\ell$, then $\ell^2$ is coprime with
\begin{qoptions}
\opt{A}{$\Bigl(1+\tan\dfrac\pi8\Bigr)^2$}
\opt{B}{$\Bigl(\tan\dfrac{3\pi}8-1\Bigr)^2$}
\opt{C}{$(\tan15^\circ-2)^2+1$}
\opt{D}{$8\sin10^\circ\sin50^\circ\sin70^\circ$}
\end{qoptions}
\qn{3} If $f(k)=\displaystyle\sum_{r=1}^{k}\dfrac1r$ and $\displaystyle\sum_{r=1}^{2023}f(r)=af(b)+c$ and $b=2023$, $a,c\in I$ then
\begin{qoptions}
\opt{A}{$b=c$}
\opt{B}{$a-b=1$}
\opt{C}{$b+c=0$}
\opt{D}{$a+c=0$}
\end{qoptions}
\qn{4} If $\alpha$ is one of the solutions of the equation $[\log_{10}4x]=[\log_{10}x]$ such that $\log_{10}\alpha<0$, then (where $[.]$ denotes greatest integer function)
\begin{qoptions}
\opt{A}{Least value of $\dfrac1\alpha$ is 4}
\opt{B}{Least integral value of $\dfrac1\alpha$ is 5}
\opt{C}{Greatest value of $\dfrac1\alpha$ is 10}
\opt{D}{$\dfrac1\alpha$ cannot be equal to 11}
\end{qoptions}
\qn{5} Let $S_n=\dfrac43+\dfrac{4.2^4}{15}+\dfrac{4.3^4}{35}+\dfrac{4.4^4}{63}+\cdots$upto $n$ terms. Then
\begin{qoptions}
\opt{A}{$n^{\mathrm{th}}$ term of sequence is $\dfrac{4n^4}{(2n+1)(2n+3)}$}
\opt{B}{$n^{\mathrm{th}}$ term of sequence is $\dfrac{4n^4}{(2n+1)(2n-1)}$}
\opt{C}{$33\{S_{16}\}=4$ where $\{\cdot\}$ denotes fractional part}
\opt{D}{$33\{S_{16}\}=8$ where $\{\cdot\}$ denotes fractional part}
\end{qoptions}
\qn{6} Consider $f(x)=ax^2+bx+c\ \forall\ x\in R$, a quadratic expression such that for $x\in[0,2]$ range of $f(x)$ is $[0,1]$ and $-\dfrac b{2a}\notin(0,2)$\\
Then which of following are correct
\begin{qoptions}
\opt{A}{The value of $C$ may be 0 or 1}
\opt{B}{The sum of roots of $ax^2+bx+c$ is $P^2+P+3$ for $P\in R$, then $|a|$ can be 2}
\opt{C}{The value of $|2a+b|$ is $\dfrac12$}
\opt{D}{If sum of roots of $ax^2+bx+c$ is $P^2+P+3$ for $P\in R$ then $|a|\le\dfrac23$}
\end{qoptions}
\sect{SECTION-I(ii)}
\paraq{7}{8}
Given that,\quad $A=2^{\log_3e^{\log_{\sqrt e}3}}$\\
$B=\sqrt{17+12\sqrt2}-\sqrt{17-12\sqrt2}$;\\
and\quad $(\sqrt2-1)^{-C}=\dfrac4{A+B}$\\
On the basis of above information, answer the following questions :
\qn{7} For three series $\dfrac{A+B}4,\ \lambda C,\ \dfrac{B-A}4\ \ldots$, select the incorrect option
\begin{qoptions}
\opt{A}{They are in G.P. if $\lambda=1$}
\opt{B}{If they are in A.P. then $\lambda=-\sqrt2$}
\opt{C}{If $\lambda=1$ then sum of infinite terms finitely exists}
\opt{D}{If $\lambda=1$ then sum of infinite terms is not finite}
\end{qoptions}
\qn{8} If 2 lies between the roots of quadratic equation $Ax^2+\sqrt2Bx+t=0$, then $|t|$ can be
\begin{qoptions}
\opt{A}{32}
\opt{B}{31}
\opt{C}{33}
\opt{D}{30}
\end{qoptions}
\paraq{9}{10}
$f(x)$ is a polynomial of degree $n$, such that $(\alpha+1)f(\alpha)-\alpha=0\ \forall\ \alpha\in N\cup\{0\}$, $\alpha\le n$
\qn{9} For $n=10$ ; $f(11)$ is
\begin{qoptions}
\opt{A}{11}
\opt{B}{$\dfrac{10}{12}$}
\opt{C}{$\dfrac{10}{11}$}
\opt{D}{1}
\end{qoptions}
\qn{10} For $n=15$, $f(16)$ is
\begin{qoptions}
\opt{A}{15}
\opt{B}{$\dfrac{15}{17}$}
\opt{C}{$\dfrac{15}{16}$}
\opt{D}{1}
\end{qoptions}
\sect{SECTION-II}
\qn{1} Let $A$, $G^2$ and $H$ be the roots of cubic equation $x^3+px^2+qx+r=0$ which are in G.P., where $p$, $q$ are integers and $A$,$G$,$H$ are respectively AM, GM, HM of two positive numbers. If $p,q\in(-100,100)$ and $\lambda$ represents the number of all possible ordered triplets $(p,q,r)$ then $\dfrac\lambda{100}$ is
\qn{2} $\dfrac{\displaystyle\sum_{r=1}^{20}\bigl(r(r+1)\bigr)^2(-1)^r}{\Bigl(\displaystyle\sum_{r=1}^{10}(4r)\Bigr)^2}$ is equal to
\qn{3} Let $T_n=1+2+3+\ldots+n$ and $M_n=\dfrac{T_2}{T_2-1}\times\dfrac{T_3}{T_3-1}\times\dfrac{T_4}{T_4-1}\ldots\times\dfrac{T_n}{T_n-1}$ for $n=2,3,4,\ldots$. The value of $52M_{50}$ is
\qn{4} Let $a_1,a_2,\ldots,a_n$, be real numbers such that\\
$\sqrt{a_1}+\sqrt{a_2-1}+\sqrt{a_3-2}+\ldots+\sqrt{a_n-(n-1)}=\dfrac12(a_1+a_2+\ldots+a_n)-\dfrac{n(n-3)}4$\\
then find the value of $\displaystyle\sum_{i=1}^{100}a_i$
\qn{5} If $x=\log_23-4\log_{\left(\ell n\frac5\pi\right)}3$, $y=\log_{\sqrt3}\Bigl(\ell n\dfrac54\Bigr)$ and $z=\log_2\bigl(\log_{(5/4)}e\bigr)$, then find the absolute value of $(xy+2z)$
\qn{6} The minimum value of $x^2+2xy+3y^2-6x-2y$ where $x,y\in R$ is equal to $m$. Then sum of all the prime number less than $|m|$ is
\qn{7} If the equation $x^4+ax^3+2x^2+bx+1=0$; $(a,b\in R)$ has a real solution, then minimum value of $a^2+b^2$ is
\qn{8} Sum of all the real solutions of the equation $x^2-3\Bigl[\sin\Bigl(x-\dfrac\pi6\Bigr)\Bigr]=3$ is $m$. Then $\dfrac{[10m]}2$ is (where $[.]$ denotes the greatest integer function)
\papersec{2}
\partsec{1}{PHYSICS}
\sect{SECTION-I(i)}
\qn{1} An air filled capacitor plate has a plate separation of $0.885\pm0.009$ mm. The area of the plates is $0.040\pm0.001$ m$^2$. It is connected to a voltage source of voltage $10.0\pm0.1$V. Choose the correct statement(s)
\begin{qoptions}
\opt{A}{The charge has a percentage error of approximately 4.5\%.}
\opt{B}{The energy stored has a percentage error of approximately 5.5\%}
\opt{C}{The force between the plates has a percentage error of approximately 5.5\%.}
\opt{D}{If we insert a dielectric of dielectric constant $k=4$ which fills the region between them completely, the random error in energy increases.}
\end{qoptions}
\qn{2} Ram performs optical bench experiment to find out the focal length of a concave lens. For this purpose, he first takes a convex lens and measures the position of object needle and image needle. Then he combines convex lens and concave lens and again measures the object needle and image needle position. The optical bench has no bench error. The readings obtained by Ram in the two experiments are shown in the table below.
\begin{center}\small
\begin{tabular}{|l|c|c|c|}
\hline & Object needle & Lens & Image needle \\\hline
Convex lens & 37.1 cm & 50.6 cm & 77.6 cm \\\hline
\begin{tabular}[c]{@{}l@{}}Convex and\\ Concave lens\end{tabular} & 37.1 cm & 50.6 cm & 91.1 cm \\\hline
\end{tabular}
\end{center}
\begin{qoptions}
\opt{A}{The focal length of the convex lens is 9.0 cm.}
\opt{B}{The focal length of the concave lens is 81 cm in modulus.}
\opt{C}{The error in focal length of convex lens is approximately 0.1 cm.}
\opt{D}{The error in focal length of concave lens is approximately 1.0 cm.}
\end{qoptions}
\qn{3} Consider an ideal solenoid having radius of cross section $R$. There are $n$ turns/length. It carries a current of $I_0\sin(\omega t)$. Consider a cylinder which is coaxial with the solenoid and has a length $\ell$ and radius $r$ ($r<R$). (Neglect magnetic field due to induced electric field)
\begin{qoptions}
\opt{A}{The rate of energy flow through its flat face is zero.}
\opt{B}{At $t=\dfrac\pi{4\omega}$, energy is flowing outwards from the curved surface of the cylinder.}
\opt{C}{The rate of energy flow is proportional to $\mu_0n^2I_0^2r^2\omega\ell$.}
\opt{D}{At $t=\dfrac\pi{4\omega}$, rate of energy flow is maximum.}
\end{qoptions}
\qn{4} Two plane mirrors, a source $S$ of light, emitting wavelength of $\lambda_1=4000$\AA\ and $\lambda_2=5600$\AA\ and a screen are arranged as shown in figure. The angle $\theta$ shown is 0.05 radian and distance $a$ and $b$ are 1 cm and 38 cm respectively.
\qfig{p2_phy_s1_q04_1}
Choose the correct options.
\begin{qoptions}
\opt{A}{fringe width of $\lambda_1=4000$\AA\ on screen is 0.08mm}
\opt{B}{Angular fringe width of $\lambda_2=5600$\AA\ on screen is $2.8\times10^{-4}$ rad}
\opt{C}{Distance of first black line from central maxima is 0.25 mm}
\opt{D}{The distance between two dark lines nearest to central maxima is 0.5 mm.}
\end{qoptions}
\qn{5} A plane electromagnetic harmonic wave of frequency $600\times10^{12}$ Hz (green light), propagating in the positive $x$ direction in vacuum, has an electric field amplitude of 42.43 V/m. The wave is linearly polarized such that the plane of vibration of the electric field is at $45^\circ$ to the xz plane.
\begin{qoptions}
\opt{A}{The equation for $\vec E$ is $\vec E=(30\text{V/m})\sin\Bigl(2\pi\bigl(600\times10^{12}\bigr)\Bigl(t-\dfrac x{3\times10^8}\Bigr)\Bigr)(\hat j+\hat k)$}
\opt{B}{The equation for $\vec B$ is $\vec B=(10^{-7}\mathrm T)\sin\Bigl(2\pi\bigl(600\times10^{12}\bigr)\Bigl(t-\dfrac x{3\times10^8}\Bigr)\Bigr)(\hat j+\hat k)$}
\opt{C}{For equation for $\vec B$ is $\vec B=(10^{-7}\mathrm T)\sin\Bigl(2\pi\bigl(600\times10^{12}\bigr)\Bigl(t-\dfrac x{3\times10^8}\Bigr)\Bigr)(-\hat j+\hat k)$}
\opt{D}{For equation for $\vec S$ is $\vec S=\Bigl(\dfrac{15}\pi\text{W/m}^2\Bigr)\sin^2\Bigl(2\pi\bigl(600\times10^{12}\bigr)\Bigl(t-\dfrac x{3\times10^8}\Bigr)\Bigr)\hat i$}
\end{qoptions}
\qn{6} A thin film of magnesium fluoride ($\mu=1.38$) is used to coat a camera lens ($\mu=1.56$). The thickness of coating which will not allow any strong reflection for the wavelength 552 nm is/are :
\begin{qoptions}
\opt{A}{1000 nm}
\opt{B}{500 nm}
\opt{C}{300 nm}
\opt{D}{100 nm}
\end{qoptions}
\sect{SECTION-I(ii)}
\qn{7} During space exploration, ISRO found a new planet similar to Earth. Experiments are conducted by four astronauts using simple pendulum to measure the value of $g$ on this new planet. The observations are shown in table below.
\begin{center}\small\setlength{\tabcolsep}{3pt}
\begin{tabular}{|c|c|c|c|}
\hline Astronaut & \begin{tabular}[c]{@{}c@{}}Length of\\ pendulum\\ (cm)\end{tabular} & \begin{tabular}[c]{@{}c@{}}number of\\ oscillations\\ ($n$)\end{tabular} & \begin{tabular}[c]{@{}c@{}}Total Time\\ (sec)\end{tabular} \\\hline
(1) & 10 & 12 & 40 \\\hline
(2) & 20 & 5 & 60 \\\hline
(3) & 5 & 10 & 30 \\\hline
(4) & 10 & 8 & 20 \\\hline
\end{tabular}
\end{center}
Least count of length is 0.1 cm and least count of time is also 0.1 sec. List-I given about four astronauts and List-II gives fractional error in the measurement of $g$.
\begin{center}\small\renewcommand{\arraystretch}{1.6}
\begin{tabular}{|l@{\hspace{4pt}}l|l@{\hspace{4pt}}l|}
\hline \multicolumn{2}{|c|}{List-I} & \multicolumn{2}{c|}{List-II} \\\hline
(P) & Astronaut-1 & (1) & $\dfrac{3}{200}$ \\\hline
(Q) & Astronaut-2 & (2) & $\dfrac{1}{120}$ \\\hline
(R) & Astronaut-3 & (3) & $\dfrac{2}{75}$ \\\hline
(S) & Astronaut-4 & (4) & $\dfrac{1}{50}$ \\\hline
\end{tabular}
\end{center}
\begin{qoptions}
\opt{A}{P $\to$ 1; Q $\to$ 3; R $\to$ 2; S $\to$ 4}
\opt{B}{P $\to$ 2; Q $\to$ 1; R $\to$ 3; S $\to$ 4}
\opt{C}{P $\to$ 3; Q $\to$ 4; R $\to$ 2; S $\to$ 1}
\opt{D}{P $\to$ 1; Q $\to$ 2; R $\to$ 3; S $\to$ 4}
\end{qoptions}
\qn{8} Two plane wavefronts $S_1=(2\text{mm})\sin(340\pi t-\pi x)$ and $S_2=2\text{mm}\sin(\pi y-340\pi t)$ are travelling along $x$ and $y$ direction respectively. They represent 2 sound waves in same medium. At different point in space, comment on motion of particles of the medium and intensity.
\begin{center}\small
\begin{tabular}{|l@{\hspace{4pt}}l|l@{\hspace{4pt}}p{3.6cm}|}
\hline \multicolumn{2}{|c|}{List-I} & \multicolumn{2}{c|}{List-II} \\\hline
(P) & (0,0) & (1) & Particles will execute SHM \\\hline
(Q) & (2,1) & (2) & Particles will undergo circular motion \\\hline
(R) & (0.5,2) & (3) & Intensity will be maximum \\\hline
(S) & (3,5) & (4) & Intensity will be minimum \\\hline
\end{tabular}
\end{center}
\begin{qoptions}
\opt{A}{P $\to$ 1,3; Q $\to$ 1,4; R $\to$ 2; S $\to$ 1,3}
\opt{B}{P $\to$ 1,3; Q $\to$ 1,4; R $\to$ 3; S $\to$ 2,4}
\opt{C}{P $\to$ 2,4; Q $\to$ 1; R $\to$ 1,4; S $\to$ 1,2}
\opt{D}{P $\to$ 3; Q $\to$ 2,4; R $\to$ 1,3; S $\to$ 3,4}
\end{qoptions}
\qn{9} Answer the following by appropriately matching the lists based on the information given in the paragraph.\\
A point source $S$ emitting monochromatic light having wavelength $\lambda$ is placed at a very small height $h=2\lambda$ above flat reflecting surface $AB$ as shown in diagram. The intensity of the reflected light is 36\% of incident intensity. Interference fringes are observed on a screen placed parallel to the reflecting surface at a large distance $D$ ($\gg h$) from it. The nature of interference fringes formed on screen is circular.
\qfig{p2_phy_s2_q09_1}
List I gives some physical quantities and List II gives some possible values of these physical quantities ($I_0$ is incident intensity)
\begin{center}\small\renewcommand{\arraystretch}{1.6}
\begin{tabular}{|l@{\hspace{4pt}}l|l@{\hspace{4pt}}l|}
\hline \multicolumn{2}{|c|}{List-I} & \multicolumn{2}{c|}{List-II} \\\hline
(P) & Intensity of maxima & (1) & $\sqrt{143}\,D$ \\\hline
(Q) & Intensity of minima & (2) & $2.56\,I_0$ \\\hline
(R) & Radius of first maxima & (3) & $0.16\,I_0$ \\\hline
(S) & Radius of second maxima & (4) & $\sqrt{63}\,D$ \\\hline
\multicolumn{2}{|c|}{} & (5) & $\dfrac{\sqrt{55}\,D}{3}$ \\\hline
\multicolumn{2}{|c|}{} & (6) & $\sqrt{15}\,D$ \\\hline
\end{tabular}
\end{center}
\begin{qoptions}
\opt{A}{P $\to$ 2; Q $\to$ 3; R $\to$ 5; S $\to$ 6}
\opt{B}{P $\to$ 2; Q $\to$ 3; R $\to$ 4; S $\to$ 5}
\opt{C}{P $\to$ 2; Q $\to$ 3; R $\to$ 5; S $\to$ 4}
\opt{D}{P $\to$ 3; Q $\to$ 2; R $\to$ 5; S $\to$ 6}
\end{qoptions}
\qn{10} A standing wave, formed in an organ pipe aligned along the x-axis, is given by\\
$S=(5\times10^{-6}\text{ m})\sin(4\pi x)\cos(1320\pi t)$\\
Where $x$ is in m and $t$ is in seconds.
\begin{center}\small\renewcommand{\arraystretch}{1.6}
\begin{tabular}{|l@{\hspace{4pt}}p{3.6cm}|l@{\hspace{4pt}}l|}
\hline \multicolumn{2}{|c|}{List-I} & \multicolumn{2}{c|}{List-II} \\\hline
(P) & Magnitude of change in pressure is maximum & (1) & $x=\dfrac14$ m \\\hline
(Q) & Magnitude of change in pressure is minimum & (2) & $x=\dfrac12$ m \\\hline
(R) & Maximum amplitude of vibration & (3) & $x=1$ m \\\hline
(S) & Minimum amplitude of vibration & (4) & $x=\dfrac18$ m \\\hline
\multicolumn{2}{|c|}{} & (5) & $x=\dfrac38$ m \\\hline
\end{tabular}
\end{center}
\begin{qoptions}
\opt{A}{P $\to$ 1, 2, 3; Q $\to$ 4, 5; R $\to$ 4, 5; S $\to$ 1, 2, 3}
\opt{B}{P $\to$ 1, 3; Q $\to$ 3, 5; R $\to$ 2, 3, 4; S $\to$ 1, 2, 3, 4, 5}
\opt{C}{P $\to$ 2, 3; Q $\to$ 2, 4; R $\to$ 2, 3, 5; S $\to$ 3, 4, 5}
\opt{D}{P $\to$ 1, 3, 4; Q $\to$ 1, 3; R $\to$ 2, 3; S $\to$ 2, 3, 4, 5}
\end{qoptions}
\sect{SECTION-II}
\qn{1} In an experiment to measure the Planck constant, $h$, the following results were obtained, in units of $10^{-34}$ Js : 6.65, 6.67. 6.57, 6.61, 6.72, 6.60, 6.71, 6.66, 6.64, 6.67. The percentage error in measurement of Planck constant is
\qn{2} A screw gauge is used to measure the thickness of a thin sheet of copper. The pitch of screw gauge is 0.2 mm and total number of division on circular scale is 200. When two jaws are brought in contact then $160^{\mathrm{th}}$ division on circular scale is exactly coinciding with the main scale line, and that the zero of main scale is not visible. When thickness of sheet is measured with screw gauge then main scale reading is 0.6 mm and $100^{\mathrm{th}}$ division coincide with the main scale line. Find the thickness of copper sheet (in mm).
\qn{3} The plane flies horizontally at an altitude of $H=4$ km at supersonic speed. The sound reaches an observer on ground 10 seconds after the plane crosses him from above. The speed of sound is 320 m/s. The speed of plane (in m/s) is
\qn{4} $S_1$ and $S_2$ are identical sources of sound of a single frequency. They are equidistant from $O$. When $S_1$ alone is sounding, the amplitude reaching $O$ is $A$. What is the increase in acoustic intensity in decibels when $S_2$ is turned on ?
\qn{5} A radar antenna is located atop a high cliff on the edge of a lake. This antenna operates on a wavelength of 400 m. Venus rises above the horizon and is tracked by the antenna. The first minimum in the signal reflected off the surface of Venus is recorded when Venus is $35^\circ$ above the horizon. The height of cliff (in m) is : ($\sin35^\circ=0.58$, $\sin20^\circ=0.34$)
\qn{6} Light of wavelength 400 nm is incident normally on a plano-convex glass lens placed on a horizontal glass surface. The radius of curvature of the curved surface is 5m. The curved face touches glass surface. The radius of third bright fringe (in mm) is
\qn{7} Polarized light of intensity $I_0$ is incident on a pair of Polaroid sheets. Let $\theta_1$ and $\theta_2$ be the angle between the incident amplitude and the axes of the first and second sheets, respectively. If $\theta_1=30^\circ$ and $\theta_2=60^\circ$, the intensity of transmitted light is $\alpha I_0$, the value of $\alpha$ is
\qn{8} In single slit fraunhoffer diffraction pattern having a slit width $a=\lambda$, resultant amplitude at the centre of central maxima is $A_0$ and it is $A$ at an angular position of $30^\circ$. Then $\dfrac{A}{A_0}$ will be
\partsec{2}{CHEMISTRY}
\sect{SECTION-I(i)}
\qn{1} Which of the following is/are correct statement/s about alkali metals.
\begin{qoptions}
\opt{A}{Reaction with water becomes increasely violent on descending the group}
\opt{B}{The reaction M$^{\oplus}$ + e$^{\ominus}$ $\to$ M has largest negative E$^\circ$ value for lithium and hence the largest positive $\Delta$G$^\circ$ value among alkali metals}
\opt{C}{KO$_2$ is used in space capsules, submarines and breathing masks because it produces dioxygen and remove carbon dioxide}
\opt{D}{All carbonate salt of alkali metal on heating give CO$_2$}
\end{qoptions}
\qn{2} In which of the following reaction/s hydrogen obtain as a product.
\begin{qoptions}
\opt{A}{Na[BH$_4$] + I$_2$ $\to$}
\opt{B}{BF$_3$ + LiAlH$_4$ $\to$}
\opt{C}{BF$_3$ + NaH $\to$}
\opt{D}{B$_2$H$_6$ + H$_2$O $\to$}
\end{qoptions}
\qn{3} Which of the following is/are correct statements about allotropes of carbon.
\begin{qoptions}
\opt{A}{C$_{60}$ contain total twenty 5-membered ring and it contain aromatic character}
\opt{B}{Graphite is composed of flat two dimentional sheets of carbon atoms each sheet is a hexagonal net of C atoms}
\opt{C}{Graphite is thermodynamically more stable than diamond}
\opt{D}{Fullerene is soluble in organic solvent}
\end{qoptions}
\qn{4} Which of the following reaction give a product which contain N-N bond.
\begin{qoptions}
\opt{A}{Mg$_3$N$_2$ $\xrightarrow{\Delta}$}
\opt{B}{NH$_4$NO$_2$ $\xrightarrow{\Delta}$}
\opt{C}{(NH$_4$)$_2$Cr$_2$O$_7$ $\xrightarrow{\Delta}$}
\opt{D}{(NH$_4$)$_2$SO$_4$ $\xrightarrow{\Delta}$}
\end{qoptions}
\qn{5} In which of the following series all the species contains bond between same atom.
\begin{qoptions}
\opt{A}{H$_2$SO$_5$, H$_2$S$_2$O$_8$, H$_2$S$_2$O$_3$}
\opt{B}{C$_3$O$_2$, S$_3$O$_9$, N$_2$O$_5$}
\opt{C}{N$_2$O, N$_2$O$_4$, N$_2$O$_3$ (unsymmetrical)}
\opt{D}{H$_2$S$_2$O$_4$, H$_2$S$_4$O$_6$, H$_2$S$_2$O$_6$}
\end{qoptions}
\qn{6} Choose correct statement about paramagnets.
\begin{qoptions}
\opt{A}{Permagnates can be prepared by treating a solution of Mn$^{2+}$ excess with very strong oxidizing agent such as PbO$_2$, or sodium bismuthate}
\opt{B}{Colour in KMnO$_4$ is due to d-d transition}
\opt{C}{When MnO$_4^{\ominus}$ react with I$^{\ominus}$ in acidic medium IO$_3^{\ominus}$ formed}
\opt{D}{KMnO$_4$ on heating does not give O$_2$}
\end{qoptions}
\sect{SECTION-I(ii)}
\qn{7}
\begin{center}\small\setlength{\tabcolsep}{3pt}
\begin{tabular}{|l@{\hspace{4pt}}l|l@{\hspace{4pt}}p{3.0cm}|}
\hline \multicolumn{2}{|c|}{\begin{tabular}[c]{@{}c@{}}\textbf{List-I}\\ \textbf{(Colour of ppt.)}\end{tabular}} & \multicolumn{2}{c|}{\begin{tabular}[c]{@{}c@{}}\textbf{List-II}\\ \textbf{(From reactions in aq. solution)}\end{tabular}} \\\hline
(P) & Red & (1) & AgNO$_3$ + NaBr $\to$ \\\hline
(Q) & Green & (2) & Hg$_2^{2+}$ + CrO$_4^{2-}$ $\to$ \\\hline
(R) & Pale yellow & (3) & Ni$^{2+}$ + NaOH $\to$ \\\hline
(S) & White & (4) & Ca$^{2+}$ + C$_2$O$_4^{2-}$ $\to$ \\\hline
\end{tabular}
\end{center}
The correct option is :-
\begin{qoptions}
\opt{A}{P $\to$ 1; Q $\to$ 4; R $\to$ 2; S $\to$ 3}
\opt{B}{P $\to$ 2; Q $\to$ 3; R $\to$ 4; S $\to$ 1}
\opt{C}{P $\to$ 2; Q $\to$ 3; R $\to$ 1; S $\to$ 4}
\opt{D}{P $\to$ 4; Q $\to$ 2; R $\to$ 1; S $\to$ 3}
\end{qoptions}
\qn{8}
\begin{center}\small\setlength{\tabcolsep}{3pt}
\begin{tabular}{|l@{\hspace{4pt}}l|l@{\hspace{4pt}}p{3.2cm}|}
\hline \multicolumn{2}{|c|}{\begin{tabular}[c]{@{}c@{}}\textbf{List-I}\\ \textbf{(Compounds)}\end{tabular}} & \multicolumn{2}{c|}{\textbf{List-II (Reagent)}} \\\hline
(P) & FeCl$_3$ & (1) & Insoluble in water as well as excess NH$_3$ solution \\\hline
(Q) & AgI & (2) & Give orange ppt with K$_3$[Fe(CN)$_6$] \\\hline
(R) & AgNO$_3$ & (3) & Give brown solution with K$_3$[Fe(CN)$_6$] \\\hline
(S) & CaCl$_2$ & (4) & Give white ppt. with K$_4$[Fe(CN)$_6$] \\\hline
\end{tabular}
\end{center}
The correct option is :-
\begin{qoptions}
\opt{A}{P $\to$ 4; Q $\to$ 3; R $\to$ 1; S $\to$ 2}
\opt{B}{P $\to$ 3; Q $\to$ 1; R $\to$ 2; S $\to$ 4}
\opt{C}{P $\to$ 3; Q $\to$ 1; R $\to$ 4; S $\to$ 2}
\opt{D}{P $\to$ 1; Q $\to$ 3; R $\to$ 4; S $\to$ 2}
\end{qoptions}
\qn{9}
\begin{center}\small\setlength{\tabcolsep}{3pt}
\begin{tabular}{|l@{\hspace{4pt}}l|l@{\hspace{4pt}}p{3.2cm}|}
\hline \multicolumn{2}{|c|}{\begin{tabular}[c]{@{}c@{}}\textbf{List-I}\\ \textbf{(Cation in aq. solution)}\end{tabular}} & \multicolumn{2}{c|}{\textbf{List-II}} \\\hline
(P) & Cr$^{3+}$ & (1) & Give white ppt with dil. HCl which soluble in hot water \\\hline
(Q) & Co$^{2+}$ & (2) & Give black ppt with H$_2$S and blue ppt with NaOH \\\hline
(R) & Pb$^{2+}$ & (3) & Give blue colour on adding conc. HCl \\\hline
(S) & Cu$^{2+}$ & (4) & Give yellow solution with NaOH + H$_2$O$_2$ \\\hline
\end{tabular}
\end{center}
The correct option is :-
\begin{qoptions}
\opt{A}{P $\to$ 4; Q $\to$ 3; R $\to$ 1; S $\to$ 2}
\opt{B}{P $\to$ 4; Q $\to$ 1; R $\to$ 2; S $\to$ 3}
\opt{C}{P $\to$ 3; Q $\to$ 1; R $\to$ 2; S $\to$ 4}
\opt{D}{P $\to$ 1; Q $\to$ 4; R $\to$ 3; S $\to$ 2}
\end{qoptions}
\qn{10}
\begin{center}\small\setlength{\tabcolsep}{3pt}
\begin{tabular}{|l@{\hspace{4pt}}l|l@{\hspace{4pt}}p{3.0cm}|}
\hline \multicolumn{2}{|c|}{\begin{tabular}[c]{@{}c@{}}\textbf{List-I}\\ \textbf{(Cations)}\end{tabular}} & \multicolumn{2}{c|}{\begin{tabular}[c]{@{}c@{}}\textbf{List-II}\\ \textbf{(Bead colour in oxidizing}\\ \textbf{flame in cold condition)}\end{tabular}} \\\hline
(P) & Cu$^{2+}$ & (1) & Brown \\\hline
(Q) & Cr$^{3+}$ & (2) & Blue \\\hline
(R) & Mn$^{2+}$ & (3) & Green \\\hline
(S) & Ni$^{2+}$ & (4) & Violet \\\hline
\end{tabular}
\end{center}
The correct option is :-
\begin{qoptions}
\opt{A}{P $\to$ 3; Q $\to$ 2; R $\to$ 1; S $\to$ 4}
\opt{B}{P $\to$ 2; Q $\to$ 3; R $\to$ 1; S $\to$ 4}
\opt{C}{P $\to$ 4; Q $\to$ 3; R $\to$ 2; S $\to$ 1}
\opt{D}{P $\to$ 2; Q $\to$ 3; R $\to$ 4; S $\to$ 1}
\end{qoptions}
\sect{SECTION-II}
\qn{1} Consider following cations in aqueous solution :-\\
Pb$^{2+}$, Hg$_2^{2+}$, Ag$^+$, Cr$^{3+}$, Al$^{3+}$, Zn$^{2+}$, Mn$^{2+}$, Co$^{2+}$\\
Calculate value of $\Bigl(\dfrac{x+y}z\Bigr)$.\\
where\\
$x$ = Total number of cation/s which give white precipitate of sulphide with Na$_2$S$_{\mathrm{(aq.)}}$\\
$y$ = Total number of cation/s which give white precipitate with dil. HCl\\
$z$ = Total number of cation/s which gives clear solution in excess NaOH solution
\qn{2} Consider the following reaction/s in aqueous solution.
\qfig{p2_chem_s3_q02_1}
if total number of unpaired electron present in complex compound (D) is $x$\\
if total number of geometrical isomer show by compound (E) than calculate $x+y$.
\qn{3} Total number of nitrate salt which give paramagnetic gas on thermal decomposition\\
LiNO$_3$, NaNO$_3$, KNO$_3$, RbNO$_3$, Pb(NO$_3$)$_2$, Mg(NO$_3$)$_2$, NH$_4$NO$_3$
\qn{4} Total number of reaction/s in which green ppt. obtain in aqueous solution.\\
(a) Ni$^{2+}$ + DMG$^{-1}$ $\to$\\
(b) CuSO$_4$ + [Fe(CN)$_6$]$^{3-}$ $\to$\\
(c) Hg$_2^{2+}$ + I$^-$ $\to$\\
(d) Ni$^{2+}$ + I$^-$ $\to$\\
(e) Ag$^+$ + PO$_4^{3-}$ $\to$\\
(f) Ag$^+$ + CO$_3^{2-}$ $\to$\\
(g) Fe$^{3+}$ + CH$_3$COO$^{\ominus}$ $\to$\\
(h) Fe$^{3+}$ + S$^{2\ominus}$ $\to$
\qn{5} Which of the following reagent/s give N$_2$ gas with NO$_2^{\ominus}$$_{\mathrm{(aq.)}}$\\
Sulphamic acid, C$_2$H$_5$NH$_2$, heating with NH$_4$Cl, Urea, Thiourea, Zn + NaOH
\qn{6} Total number of elements/ions which give H$_2$ with dilute HCl.\\
Na, Zn, Cu, Ag, Cr$^{2+}$, Mn$^{2+}$, Ti$^{2+}$, Fe$^{2+}$, Co$^{2+}$
\qn{7} 1. XeF$_4$ + SbF$_5$ $\to$ product (P)\\
2. XeF$_4$ + O$_2$F$_2$ $\to$ product (Q) + paramagnetic gas (G)\\
Consider above reaction and select total number of correct statement/s.\\
(A) In product (Q) formation O$_2$F$_2$ act as Fluorinating agent\\
(B) Product (Q) in solid state contain distorted octahedral structure\\
(C) When product (Q) undergoes complete hydrolysis than it give explosive compound\\
(D) Product (P) is ionic compound and Xe atom is anionic part\\
(E) Product (P) is ionic compound and Xe atom in cationic part\\
(F) In product (P) no ionic product is formed\\
(G) In product (P) Xe atom has sp$^3$d$^2$ hybridization
\qn{8} Consider following compound/s :-\\
XeF$_2$, XeF$_4$, XeF$_6$, XeOF$_4$, XeO$_3$, XeO$_2$F$_2$\\
if (A) is/are total number of compound/s which give oxygen when react with water\\
(B) is/are total number of compound/s which show non-redox reaction when added in water\\
(C) is/are total number of planar compound/s\\
then calculate $\dfrac{(A+B)}C$
\partsec{3}{MATHEMATICS}
\sect{SECTION-I(i)}
\qn{1} If the real values of `$a$' for which the equation $\bigl(x-\log_2(a^6+3a^2+3a^4+1)\bigr)(x-1)+1=0$ has integral roots are $a_1$ and $a_2$, then
\begin{qoptions}
\opt{A}{$a_1+a_2=0$}
\opt{B}{$a_1^2+a_2^2=2$}
\opt{C}{$|a_1-a_2|=2$}
\opt{D}{$|a_1-a_2|=0$}
\end{qoptions}
\qn{2} Possible set of values of `$a$' for which the equation $(x^2+x)^2+a(x^2+x)+4=0$ have all four imaginary roots is/are
\begin{qoptions}
\opt{A}{$\Bigl(\dfrac{65}4,\infty\Bigr)$}
\opt{B}{$\Bigl(-4,\dfrac{65}4\Bigr)$}
\opt{C}{$a\in\Bigl(-4,\dfrac{65}8\Bigr)$}
\opt{D}{$\Bigl(\dfrac{65}8,\infty\Bigr)$}
\end{qoptions}
\qn{3} If $a,b\in R$ satisfy the equation $a^3+24ab+64=8b^3$, then roots $\alpha,\beta$ ($\alpha<\beta$) of equation $x^2+bx+a=0$ satisfies (provided $0<b<3$)
\begin{qoptions}
\opt{A}{$\alpha\in(-4,-2)$}
\opt{B}{$\beta\in(-1,2)$}
\opt{C}{$\alpha\in(-4,-2]$}
\opt{D}{$\beta\in(-4,-2)$}
\end{qoptions}
\qn{4} If equation $\Bigl(\log_{3a}(a^2-4)\Bigr)x^2+bx+1=0$ doesn't have any real roots (provided all coefficient are real), then point possible set of values of $a$ contains
\begin{qoptions}
\opt{A}{$(-2,-1)$}
\opt{B}{$(1,2)$}
\opt{C}{$(3,2\sqrt5)$}
\opt{D}{$(\sqrt5,\infty)$}
\end{qoptions}
\qn{5} If the equation $e^{4x}-ae^{3x}+11e^{2x}-ae^x+1=0$ has two distinct positive roots then $a$ may take the value/s
\begin{qoptions}
\opt{A}{6}
\opt{B}{$\dfrac{13}2$}
\opt{C}{$\dfrac{25}4$}
\opt{D}{$\dfrac{19}3$}
\end{qoptions}
\qn{6} If $S_n$ denotes the sum to $n$ terms of series ($1\le n\le9$)\\
$1+22+333+\ldots+\underbrace{99\ldots9}_{9\text{ times}}$\\
then for $n\ge2$
\begin{qoptions}
\opt{A}{$S_n-S_{n-1}=\dfrac19\bigl(10^n-n^2+n\bigr)$}
\opt{B}{$S_n=\dfrac19\bigl(10^n-n^2+2n-2\bigr)$}
\opt{C}{$9(S_n-S_{n-1})=n(10^n-1)$}
\opt{D}{$S_3=356$}
\end{qoptions}
\sect{SECTION-I(ii)}
\qn{7} If $\log_32=a$, $\log_45=b$ and $\log_56=c$ then match the following columns
\begin{center}\small\renewcommand{\arraystretch}{1.6}\setlength{\tabcolsep}{3pt}
\begin{tabular}{|l@{\hspace{4pt}}l|l@{\hspace{4pt}}l|}
\hline \multicolumn{2}{|c|}{\textbf{List-I}} & \multicolumn{2}{c|}{\textbf{List-II}} \\\hline
(P) & $\log_{10}135$ & (1) & $\dfrac{3+2ab}{a(1+2b)}$ \\\hline
(Q) & $\log_436$ & (2) & $\dfrac{2bc+1}{2bc+2}$ \\\hline
(R) & $\log_{24}12$ & (3) & $\dfrac{a+1}a$ \\\hline
(S) & $2\log_{10}3+\dfrac{1+\log_53}{1+\log_52}$ & (4) & $2bc$ \\\hline
\end{tabular}
\end{center}
\begin{qoptions}
\opt{A}{P $\to$ 1; Q $\to$ 3,4; R $\to$ 2; S $\to$ 1}
\opt{B}{P $\to$ 3,4; Q $\to$ 1; R $\to$ 1; S $\to$ 2}
\opt{C}{P $\to$ 2; Q $\to$ 3,4; R $\to$ 4; S $\to$ 1}
\opt{D}{P $\to$ 4; Q $\to$ 2; R $\to$ 3; S $\to$ 3,4}
\end{qoptions}
\qn{8} If $x,y,z,u>0$ are distinct real numbers in G.P. and $a,b,c,d$ are real \& distinct, non-zero numbers in A.P., then match the correct answers.
\begin{center}\small\renewcommand{\arraystretch}{1.6}\setlength{\tabcolsep}{3pt}
\begin{tabular}{|l@{\hspace{4pt}}p{3.3cm}|l@{\hspace{4pt}}p{2.4cm}|}
\hline \multicolumn{2}{|c|}{\textbf{List-I}} & \multicolumn{2}{c|}{\textbf{List-II}} \\\hline
(P) & $a-\log\dfrac x{y^3},\ b-\log\dfrac y{z^3},\ c-\log\dfrac z{u^3}$ will be & (1) & only in A.P. \\\hline
(Q) & $(b-a)xy,\ (c-b)yz,\ (d-c)zu$ will be & (2) & only in G.P. \\\hline
(R) & $\dfrac{a(b+c)}{bc},\ \dfrac{b(c+a)}{ca},$\newline $\dfrac{c(a+b)}{ab}$ will be & (3) & only in H.P. \\\hline
(S) & $(a^2-b^2)(x^2-y^2),\ (b^2-c^2)(y^2-z^2),\ (c^2-d^2)(z^2-u^2)$ will be & (4) & only in A.G.P. \\\hline
\multicolumn{2}{|c|}{} & (5) & none of above \\\hline
\end{tabular}
\end{center}
\begin{qoptions}
\opt{A}{P $\to$ 2; Q $\to$ 1; R $\to$ 3; S $\to$ 4}
\opt{B}{P $\to$ 1; Q $\to$ 2; R $\to$ 1; S $\to$ 4}
\opt{C}{P $\to$ 3; Q $\to$ 4; R $\to$ 1; S $\to$ 2}
\opt{D}{P $\to$ 2; Q $\to$ 4; R $\to$ 5; S $\to$ 1}
\end{qoptions}
\end{multicols}
\noindent\begin{minipage}{\textwidth}\qn{9}\par\vspace{-2pt}
\begin{center}\small\renewcommand{\arraystretch}{1.6}\setlength{\tabcolsep}{4pt}
\begin{tabular}{|l@{\hspace{4pt}}p{8.2cm}|l@{\hspace{4pt}}p{4.4cm}|}
\hline \multicolumn{2}{|c|}{\textbf{List-I}} & \multicolumn{2}{c|}{\textbf{List-II}} \\\hline
(P) & If $p^{\mathrm{th}}$, $q^{\mathrm{th}}$ and $r^{\mathrm{th}}$ terms of a G.P. (common ratio $\ne1$) are in G.P. then $p,q,r$ will be in & (1) & Arithmetic progression \\\hline
(Q) & If $a,b,c,d$ are non-zero real numbers such that $(a^2+b^2+c^2)(b^2+c^2+d^2)<(ab+bc+cd)^2$ and $c(a-b)=a(b-c)$, then $a,b,c,d$ will be in & (2) & Geometric progression \\\hline
(R) & If $a,b,c,d$ are four different positive numbers in G.P. then $\dfrac{a(2\log a+\log b)}b,\ \dfrac{b(2\log b+\log c)}c,\ \dfrac{c(2\log c+\log d)}d$ will be in & (3) & Harmonic progression \\\hline
(S) & If $a_1,a_2,a_3$ are in A.P., $a_2,a_3,a_4$ are in G.P. and $a_3,a_4,a_5$ are in H.P. then $a_1,a_3,a_5$ are always in & (4) & Arithmetic geometric progression \\\hline
\multicolumn{2}{|c|}{} & (5) & A progression, whose all the terms are identical \\\hline
\end{tabular}
\end{center}\end{minipage}
\begin{multicols}{2}
\begin{qoptions}
\opt{A}{P $\to$ 2,5; Q $\to$ 3,4; R $\to$ 3,5; S $\to$ 1,5}
\opt{B}{P $\to$ 1,5; Q $\to$ 3,5; R $\to$ 3,4; S $\to$ 2,5}
\opt{C}{P $\to$ 1,4; Q $\to$ 1,2,3,4,5; R $\to$ 1,4; S $\to$ 2,4}
\opt{D}{P $\to$ 2,4; Q $\to$ 1,4; R $\to$ 2,3; S $\to$ 1,5}
\end{qoptions}
\end{multicols}
\noindent\begin{minipage}{\textwidth}\qn{10}\par\vspace{-2pt}
\begin{center}\small\renewcommand{\arraystretch}{1.6}\setlength{\tabcolsep}{4pt}
\begin{tabular}{|l@{\hspace{4pt}}p{8.2cm}|l@{\hspace{4pt}}p{4.4cm}|}
\hline \multicolumn{2}{|c|}{\textbf{List-I}} & \multicolumn{2}{c|}{\textbf{List-II}} \\\hline
(P) & If $A=\Bigl\{x\Bigm|\log_{\frac23(|x-2|)}2\ge0;\ x\in R\Bigr\}$, then possible subset of $A$ is/are & (1) & $\Bigl(0,\dfrac13\Bigr)$ \\\hline
(Q) & Subset of interval in which `$a$' lies if $x^2-(\alpha-\beta)x+(1-\alpha-\beta)=0$ has distinct real roots for all real $\beta$, is & (2) & $(1,2)$ \\\hline
(R) & The interval in which `$b$' lies, if the roots of $x^3-6x^2+2ax-4b=0$, where $a\in(5,6)$, are in A.P., is & (3) & $\Bigl(2,\dfrac72\Bigr)$ \\\hline
(S) & Subset of set of values of $\lambda$ for which both roots of equation $x^2-6\lambda x+2-2\lambda+9\lambda^2=0$ are greater than 3, is & (4) & $(-\infty,1)\cup\Bigl(\dfrac{11}9,\infty\Bigr)$ \\\hline
\end{tabular}
\end{center}\end{minipage}
\begin{multicols}{2}
\begin{qoptions}
\opt{A}{P $\to$ 1; Q $\to$ 2,3; R $\to$ 2; S $\to$ 4}
\opt{B}{P $\to$ 4; Q $\to$ 2; R $\to$ 3; S $\to$ 1}
\opt{C}{P $\to$ 3; Q $\to$ 4; R $\to$ 1; S $\to$ 2}
\opt{D}{P $\to$ 2; Q $\to$ 1; R $\to$ 4; S $\to$ 3}
\end{qoptions}
\sect{SECTION-II}
\qn{1} Find the sum of infinite terms of a decreasing G.P. whose first term is equal to $p+2$ and the common ratio is $\dfrac2q$, where $p$ is the least value of the product of the roots of the equation $(m^2+1)x^2-3x+(m^2+1)^2=0$ and $q$ is the greatest value of the sum of its roots.
\qn{2} If the sum of the series $\dfrac x{1-x^2}+\dfrac{x^2}{1-x^4}+\dfrac{x^4}{1-x^8}+\ldots$ upto 10 terms is $\dfrac1{1-x}-\dfrac1{1-x^k}$, then the value of $k$ must be
\qn{3} Evaluate the sum $\displaystyle\sum_{k=1}^{\infty}\dfrac{6^k}{(3^k-2^k)(3^{k+1}-2^{k+1})}$.
\qn{4} If the range of values of $k$ for which there is at least one common solution of the inequalities $x^2+4kx+3k^2>1+2k$ and $x^2+2kx\le3k^2-8k+4$, is $(-\infty,a)\cup(b,\infty)$, then find the value of $(a+b)$.
\qn{5} When $x^{100}$ is divided by $x^2-3x+2$, then remainder is $(2^{k+1}-1)x-2(2^k-1)$ where $k$ is a numerical quantity. Then $k$ must be
\qn{6} Given that $\alpha,\gamma$ are roots of the equation $Ax^2-4x+1=0$ and $\beta,\delta$ the roots of the equation $Bx^2-6x+1=0$. If the value of $A$ and $B$ such that $\alpha,\beta,\gamma$ and $\delta$ are in H.P. are $\lambda$ and $\lambda+5$ respectively, then the numerical quantity $\lambda$ must be
\qn{7} If $a_1,a_2,a_3\ldots a_{100}$ are in H.P. then $\displaystyle\sum_{i=1}^{99}\dfrac{a_ia_{i+1}}{a_1a_{100}}$ equals
\qn{8} If $a,b,c$ are positive and $a+b+c=1$, then find minimum value of $\Bigl(\dfrac1a-1\Bigr)\Bigl(\dfrac1b-1\Bigr)\Bigl(\dfrac1c-1\Bigr)$ equals

\end{multicols}
\end{document}
